218
13 Line and Surface Integrals
It remains only to recall that the coordinate representation $f(t) dt^1 \wedge dt^2$ of the
form $\varphi^*\omega^2$ is obtained from the coordinate expression for the form $\omega^2$ by the direct
substitution $x = \varphi(t)$, where $\varphi : I \to G$ is a chart of the surface $S$.
Carrying out this change of variable, we obtain from (13.4)
$$ F = \int_{S=\varphi(I)} \omega^2 = \int_I \varphi^*\omega^2 = \int_I \left[ V^1(\varphi(t)) \frac{\partial x^2}{\partial t^1} \frac{\partial x^3}{\partial t^2} - \frac{\partial x^3}{\partial t^1} \frac{\partial x^2}{\partial t^2} + V^2(\varphi(t)) \frac{\partial x^3}{\partial t^1} \frac{\partial x^1}{\partial t^2} - \frac{\partial x^1}{\partial t^1} \frac{\partial x^3}{\partial t^2} \right] dt^1 \wedge dt^2 $$
$$ \left. + V^3(\varphi(t)) \frac{\partial x^1}{\partial t^1} \frac{\partial x^2}{\partial t^2} - \frac{\partial x^2}{\partial t^1} \frac{\partial x^1}{\partial t^2} \right] dt^1 dt^2. $$
This last integral, as Eq. (13.5) shows, is the ordinary Riemann integral over the
rectangle $I$.
Thus we have found that
$$ F = \int_I \begin{pmatrix} V^1(\varphi(t)) \\ V^2(\varphi(t)) \\ V^3(\varphi(t)) \end{pmatrix} \cdot \begin{pmatrix} \frac{\partial \varphi^1}{\partial t^1} & \frac{\partial \varphi^1}{\partial t^2} \\ \frac{\partial \varphi^2}{\partial t^1} & \frac{\partial \varphi^2}{\partial t^2} \\ \frac{\partial \varphi^3}{\partial t^1} & \frac{\partial \varphi^3}{\partial t^2} \end{pmatrix} \begin{pmatrix} dt^1 \\ dt^2 \end{pmatrix} dt^1 dt^2, $$
$$ F = \int_I \begin{pmatrix} V^1(\varphi(t)) \\ V^2(\varphi(t)) \\ V^3(\varphi(t)) \end{pmatrix} \cdot \begin{pmatrix} \frac{\partial \varphi^1}{\partial t^1} & \frac{\partial \varphi^2}{\partial t^1} & \frac{\partial \varphi^3}{\partial t^1} \\ \frac{\partial \varphi^1}{\partial t^2} & \frac{\partial \varphi^2}{\partial t^2} & \frac{\partial \varphi^3}{\partial t^2} \end{pmatrix} \begin{pmatrix} dt^1 \\ dt^2 \end{pmatrix} dt^1 dt^2, \quad \text{(13.6)} $$
where $x = \varphi(t) = (\varphi^1, \varphi^2, \varphi^3)(t_1, t_2)$ is a chart of the surface $S$ defining the same
orientation as the field of normals we have given. If the chart $\varphi : I \to S$ gives $S$ the
opposite orientation, Eq. (13.6) does not generally hold. But, as follows from the
considerations at the beginning of this subsection, the left- and right-hand sides will
differ only in sign in that case.
The final formula (13.6) is obviously merely the limit of the sums of the elemen-
tary fluxes $\Delta F_i \approx (V(x_i), \xi_1, \xi_2)$ familiar to us, written accurately in the coordi-
nates $t^1$ and $t^2$.
We have considered the case of a surface defined by a single chart. In general
a smooth surface can be decomposed into smooth pieces $S_i$ having essentially no
intersections with one another, and then we can find the flux through $S$ as the sum
of the fluxes though the pieces $S_i$.
Example 3 Suppose a medium is advancing with constant velocity $V = (1, 0, 0)$.
If we take any closed surface in the domain of the flow, then, since the density of
the medium does not change, the amount of matter in the volume bounded by this
surface must remain constant. Hence the total flux of the medium through such a
surface must be zero.
In this case, let us check formula (13.6) by taking $S$ to be the sphere $x^2 + y^2 +
z^2 = R^2$.